\scanpage{62}
% archive sheet number: 61
est exacte; donc, comme $\underline{O}_X$ est plat sur $Y$, il s'ensuit (réf) que $\underline{O}_{X'}$ est plat sur $Y$, de plus $\underline{O}_{X'}\otimes k(z)$ est \inserted{\unsure{un quotient|en général}} divisoriel pour $\underline{O}_X\otimes k(z)$. \inserted{Faisons donc $z=y$.}

\uline{Suffisance} La suite exacte $0\to\underline{\unsure{J|\mathcal{J}|I}}\to\underline{O}_X\unsure{-|\to}\underline{O}_{X'}-0$ donne, puisque $\underline{O}_{X'}$ est $Y$-plat \uline{en $x$}, la suite exacte
\[
0\to\underline{J}\otimes k(y)\to\underline{O}_X\otimes k(y)\to\underline{O}_{X'}\otimes k(y)-0,
\]
et comme $\underline{O}_{X'}\otimes k(y)$ est divisoriel \inserted{en $x$}, il s'ensuit que $(\underline{J}\otimes k(y))_x$ est monogène, donc (Nakayama) $\underline{J}$ \deleted{\unsure{alors}} est \deleted{\unsure{engendré par un élém.}} \inserted{monogène} au voisinage de $x$, soit $\underline{J}=\unsure{\varphi|\phi}\underline{O}_X$ [en rapetissant $X$]. Comme $\varphi$ dans $\underline{O}_X\otimes k(y)$ est un générateur de $\underline{J}\otimes k(y)$, il n'est pas div. de $0$ dans $\underline{O}_X\otimes k(y)$, i.e. la suite
\[
0-\underline{O}_X\otimes k(y)\xrightarrow{\varphi}\underline{O}_X\otimes k(y)\to\underline{O}_{X'}\otimes k(y)\to0
\]
déduite de la suite
\[
0-\underline{O}_X\xrightarrow{\varphi}\underline{O}_X\to\underline{O}_{X'}\to0
\]
est \uline{exacte}. Je dis que $\varphi\colon\inserted{\underline{O}_X\to\underline{O}_X}$ est \uline{injectif} en $x$ : soit $I$ son noyau, on aura alors les suites exactes
\[
0-I\to\underline{O}_X\xrightarrow{\varphi}\deleted{\unsure{I}}\unsure{\underline{J}|J}\to0,\quad 0-\underline{J}\to\underline{O}_X\to\underline{O}_{X'}\to0,
\]
or $J$ est plat (car $\underline{O}_X$ et $\underline{O}_{X'}$ le sont) donc on a la suite exacte
\[
0-I\otimes k\to\underline{O}_X\otimes k\to J\otimes k\to0
\]
d'où $I\otimes k=0$ d'où $I=0$ (Nakayama). \unsure{c.q.f.d.}

\uline{Cor} Stabilité par extension de la base.
